\chapter{2015年湖南卷（文）}

\section{单选题}

\begin{problem}
已知 \(\frac{(1 - \i)^2}{z} = 1 + \i\)（\(\i\) 为虚数单位），则复数 \(z =\)（\quad）

\choices
    {\(1 + \i\)}
    {\(1 - \i\)}
    {\(-1 + \i\)}
    {\(-1 - \i\)}
\end{problem}

\begin{answer}
D
\end{answer}

\begin{solution}
因为
\[
(1-\i)^2=1-2\i+\i^2=-2\i
\]
所以
\[
z=\frac{-2\i}{1+\i}=\frac{-2\i(1-\i)}{(1+\i)(1-\i)}=-1-\i
\]
故选 D.
\end{solution}

\begin{problem}
在一次马拉松比赛中，\(35\) 名运动员的成绩（单位：\(\text{分钟}\)）茎叶图如图所示.
\renewcommand{\arraystretch}{1.2}
\begin{tabular}{r|l}
\(13\) & \(0\)，\(0\)，\(3\)，\(4\)，\(5\)，\(6\)，\(6\)，\(8\)，\(8\)，\(8\)，\(9\) \\
\(14\) & \(1\)，\(1\)，\(1\)，\(2\)，\(2\)，\(2\)，\(3\)，\(3\)，\(4\)，\(4\)，\(5\)，\(5\)，\(5\)，\(6\)，\(6\)，\(7\)，\(8\) \\
\(15\) & \(0\)，\(1\)，\(2\)，\(2\)，\(3\)，\(3\)，\(3\)
\end{tabular}

若将运动员按成绩由好到差编为 \(1\sim 35\) 号，再用系统抽样的方法从中抽取 \(7\) 人，则其中成绩在区间\([139, 151]\)上的运动员的人数是（\quad）

\choices
    {\(3\)}
    {\(4\)}
    {\(5\)}
    {\(6\)}
\end{problem}

\begin{answer}
B
\end{answer}

\begin{solution}
按成绩由好到差排列后，前 \(10\) 名的成绩为 \(130\) 至 \(138\) 分钟，第 \(11\) 名成绩为 \(139\) 分钟，第 \(12\) 至第 \(28\) 名成绩为 \(141\) 至 \(148\) 分钟，第 \(29\)、\(30\) 名成绩分别为 \(150\)、\(151\) 分钟，因此成绩在 \([139,151]\) 上的运动员对应第 \(11\) 至第 \(30\) 名.

系统抽样的抽样间隔为 \(\dfrac{35}{7}=5\). 设第一个抽取的号码为 \(r\)，则 \(r\in\{1,2,3,4,5\}\)，抽取的号码为 \(r\)，\(r+5\)，\(\ldots\)，\(r+30\). 当 \(r=1\)，\(2\)，\(3\)，\(4\)，\(5\) 时，落在第 \(11\) 至第 \(30\) 名中的号码分别为
\[
\begin{gathered}
11,16,21,26;\quad 12,17,22,27;\quad 13,18,23,28;\\
14,19,24,29;\quad 15,20,25,30.
\end{gathered}
\]
每种起始号码都抽到 \(4\) 名，所以人数为 \(4\). 故选 B.
\end{solution}

\begin{problem}
设 \(x \in \R\)，则“\(x > 1\)”是“\(x^3 > 1\)”的（\quad）

\choices
    {充分不必要条件}
    {必要不充分条件}
    {充要条件}
    {既不充分也不必要条件}
\end{problem}

\begin{answer}
C
\end{answer}

\begin{solution}
因为
\[
x^3-1=(x-1)(x^2+x+1)=(x-1)\left[\left(x+\frac12\right)^2+\frac34\right]
\]
且 \(x^2+x+1>0\)，所以 \(x^3>1\) 当且仅当 \(x-1>0\)，即 \(x>1\). 因而“\(x>1\)”是“\(x^3>1\)”的充要条件. 故选 C.
\end{solution}

\begin{problem}
若变量 \(x\)，\(y\) 满足约束条件 \(\begin{cases}
x + y \geqslant 1,\\
y - x \leqslant 1,\\
x \leqslant 1,
\end{cases}\) 则 \(z = 2x - y\) 的最小值为（\quad）

\choices
    {\(-1\)}
    {\(0\)}
    {\(1\)}
    {\(2\)}
\end{problem}

\begin{answer}
A
\end{answer}

\begin{solution}
由 \(x+y\geqslant1\) 和 \(y-x\leqslant1\)，得
\[
1-x\leqslant y\leqslant x+1
\]
要使可取的 \(y\) 存在，必须有 \(1-x\leqslant x+1\)，即 \(x\geqslant0\). 结合 \(x\leqslant1\)，可知 \(0\leqslant x\leqslant1\).

对固定的 \(x\)，\(z=2x-y\) 随 \(y\) 增大而减小，所以取 \(y=x+1\) 时 \(z\) 最小，此时
\[
z=2x-(x+1)=x-1\geqslant-1
\]
当 \(x=0\)，\(y=1\) 时取等号，故 \(z\) 的最小值为 \(-1\). 故选 A.
\end{solution}

\begin{problem}
执行如图所示的程序框图，如果输入 \(n=3\)，则输出的 \(S=\)（\quad）

\begin{minipage}{\linewidth}
    \centering
    \texfigure[max width=0.44\linewidth]{img_repaint/2015/hunan_liberal/q5_fig1.tex}
\end{minipage}

\choices
    {\(\frac{6}{7}\)}
    {\(\frac{3}{7}\)}
    {\(\frac{8}{9}\)}
    {\(\frac{4}{9}\)}
\end{problem}

\begin{answer}
B
\end{answer}

\begin{solution}
程序从 \(i=1\)、\(S=0\) 开始，每次将
\[
\frac{1}{(2i-1)(2i+1)}
\]
加入 \(S\)，再令 \(i=i+1\). 输入 \(n=3\) 时，循环在 \(i=1\)，\(2\)，\(3\) 时执行，因此
\[
S=\frac{1}{1\cdot3}+\frac{1}{3\cdot5}+\frac{1}{5\cdot7}
 =\frac13+\frac1{15}+\frac1{35}=\frac37
\]
故选 B.
\end{solution}

\begin{problem}
若双曲线 \(\frac{x^2}{a^2} - \frac{y^2}{b^2} = 1\) 的一条渐近线经过点\((3,-4)\)，则此双曲线的离心率为（\quad）

\choices
    {\( \frac{\sqrt{7}}{3} \)}
    {\( \frac{5}{4} \)}
    {\( \frac{4}{3} \)}
    {\( \frac{5}{3} \)}
\end{problem}

\begin{answer}
D
\end{answer}

\begin{solution}
双曲线的渐近线为 \(y=\pm\dfrac{b}{a}x\). 由于其中一条渐近线经过 \((3,-4)\)，且 \(a\)，\(b>0\)，所以
\[
\frac{b}{a}=\frac43
\]
双曲线的离心率满足 \(c^2=a^2+b^2\)，故
\[
e=\frac{c}{a}=\sqrt{1+\frac{b^2}{a^2}}=\sqrt{1+\frac{16}{9}}=\frac53
\]
故选 D.
\end{solution}

\begin{problem}
若实数 \(a\)，\(b\) 满足 \(\frac{1}{a} + \frac{2}{b} = \sqrt{ab}\)，则 \(ab\) 的最小值为（\quad）

\choices
    {\(\sqrt{2}\)}
    {\(2\)}
    {\(2 \sqrt{2}\)}
    {\(4\)}
\end{problem}

\begin{answer}
C
\end{answer}

\begin{solution}
由根式有意义且分母不为零，知 \(ab>0\). 若 \(a\)，\(b<0\)，则左端 \(\dfrac1a+\dfrac2b<0\)，而右端 \(\sqrt{ab}>0\)，矛盾，所以 \(a\)，\(b>0\).

令 \(t=ab\). 原等式两边同乘以 \(ab\)，得
\[
b+2a=(ab)^{3/2}=t^{3/2}
\]
由基本不等式
\[
b+2a\geqslant2\sqrt{2ab}=2\sqrt{2t}
\]
可得 \(t^{3/2}\geqslant2\sqrt{2t}\). 因为 \(t>0\)，两边除以 \(\sqrt t\)，得
\[
t\geqslant2\sqrt2
\]
当 \(b=2a\)，且 \(ab=2\sqrt2\) 时等号成立，因此 \(ab\) 的最小值为 \(2\sqrt2\). 故选 C.
\end{solution}

\begin{problem}
设函数 \( f(x) = \ln (1 + x) - \ln (1 - x) \)，则 \( f(x) \) 是（\quad）

\choices
    {奇函数，且在 \((0,1)\) 是增函数}
    {奇函数，且在 \((0,1)\) 是减函数}
    {偶函数，且在 \((0,1)\) 是增函数}
    {偶函数，且在 \((0,1)\) 是减函数}
\end{problem}

\begin{answer}
A
\end{answer}

\begin{solution}
函数定义域为 \((-1,1)\). 对任意 \(x\in(-1,1)\)，有
\[
f(-x)=\ln(1-x)-\ln(1+x)=-f(x)
\]
所以 \(f(x)\) 是奇函数. 又在 \((0,1)\) 上
\[
f'(x)=\frac{1}{1+x}+\frac{1}{1-x}=\frac{2}{1-x^2}>0
\]
故 \(f(x)\) 在 \((0,1)\) 上是增函数. 故选 A.
\end{solution}

\begin{problem}
已知点 \(A\)，\(B\)，\(C\) 在圆 \(x^{2} + y^{2} = 1\) 上运动，且 \(AB \perp BC\). 若点 \(P\) 的坐标为 \((2,0)\)，则 \(\left|\overrightarrow{PA} + \overrightarrow{PB} + \overrightarrow{PC}\right|\) 的最大值为（\quad）

\choices
    {\(6\)}
    {\(7\)}
    {\(8\)}
    {\(9\)}
\end{problem}

\begin{answer}
B
\end{answer}

\begin{solution}
设圆心为 \(O\)，以 \(O\) 为原点的向量分别为 \(\bs a\)，\(\bs b\)，\(\bs c\)，则
\[
|\bs a|=|\bs b|=|\bs c|=1
\]
由 \(AB\perp BC\)，知 \(\angle ABC=90^{\circ}\). 根据圆周角定理，\(AC\) 是圆的直径，所以
\[
\bs a+\bs c=\bs0
\]
又 \(\overrightarrow{OP}=(2,0)\)，因此
\[
\overrightarrow{PA}+\overrightarrow{PB}+\overrightarrow{PC}
 =\bs a+\bs b+\bs c-3\overrightarrow{OP}
 =\bs b-(6,0)
\]
由三角不等式
\[
\left|\bs b-(6,0)\right|\leqslant|\bs b|+|(6,0)|=1+6=7
\]
当 \(\bs b=(-1,0)\) 时取等号，例如取 \(A=(0,1)\)、\(C=(0,-1)\)，此时仍有 \(AB\perp BC\). 所以所求最大值为 \(7\). 故选 B.
\end{solution}

\begin{problem}
某工件的三视图如图所示，现将该工件通过切削，加工成一个体积尽可能大的正方体新工件，并使新工件的一个面落在原工件的一个面内，则原工件材料的利用率为（\quad）（材料的利用率 \(=\frac{\text{新工件的体积}}{\text{原工件的体积}}\)）

\begin{minipage}{\linewidth}
    \centering
    \bitmapfigure[max width=0.36\linewidth]{img/2015/hunan_liberal/q10_fig1.png}
\end{minipage}

\choices
    {\(\frac{8}{9\pi}\)}
    {\(\frac{8}{27\pi}\)}
    {\(\frac{24(\sqrt{2} - 1)^3}{\pi}\)}
    {\(\frac{8(\sqrt{2} - 1)^3}{\pi}\)}
\end{problem}

\begin{answer}
A
\end{answer}

\begin{solution}
由三视图可知，原工件是底面半径为 \(1\)、母线长为 \(3\) 的直圆锥. 设圆锥的高为 \(H\)，由勾股定理
\[
H=\sqrt{3^2-1^2}=2\sqrt2
\]
设正方体棱长为 \(s\)，且正方体的底面落在圆锥底面内. 在距圆锥底面 \(s\) 高度处，圆锥的截面半径为
\[
R=1-\frac{s}{2\sqrt2}
\]
正方形的对角线不能超过该圆的直径，所以正方体顶面落在该截面内的充要条件为
\[
\frac{s}{\sqrt2}\leqslant R=1-\frac{s}{2\sqrt2}
\]
解得 \(s\leqslant\dfrac{2\sqrt2}{3}\). 取最大值 \(s=\dfrac{2\sqrt2}{3}\)，新工件体积为
\[
V_{\text{新}}=s^3=\left(\frac{2\sqrt2}{3}\right)^3=\frac{16\sqrt2}{27}
\]
原工件体积为
\[
V_{\text{原}}=\frac13\pi\cdot1^2\cdot2\sqrt2=\frac{2\sqrt2\pi}{3}
\]
因此材料的利用率为
\[
\frac{V_{\text{新}}}{V_{\text{原}}}
 =\frac{16\sqrt2/27}{2\sqrt2\pi/3}=\frac{8}{9\pi}
\]
故选 A.
\end{solution}

\section{填空题}

\begin{problem}
已知集合 \(U = \{1, 2, 3, 4\}\)，\(A = \{1, 3\}\)，\(B = \{1, 3, 4\}\)，则 \(A \cup \complement_U B = \)\fillinblank{}.
\end{problem}

\begin{answer}
\(\{1,2,3\}\)
\end{answer}

\begin{solution}
由 \(U=\{1,2,3,4\}\)、\(B=\{1,3,4\}\)，得
\[
\complement_U B=\{2\}
\]
所以
\[
A\cup\complement_U B=\{1,3\}\cup\{2\}=\{1,2,3\}
\]
\end{solution}

\begin{problem}
在直角坐标系 \(xOy\) 中，以坐标原点为极点，\(x\) 轴的正半轴为极轴建立极坐标系. 若曲线 \(C\) 的极坐标方程为 \(\rho = 2\sin \theta\)，则曲线 \(C\) 的直角坐标方程为\fillinblank{}.
\end{problem}

\begin{answer}
\(x^2+y^2-2y=0\)
\end{answer}

\begin{solution}
由极坐标与直角坐标的关系 \(\rho^2=x^2+y^2\)、\(\rho\sin\theta=y\)，将 \(\rho=2\sin\theta\) 两边同乘以 \(\rho\)，得
\[
\rho^2=2\rho\sin\theta
\]
因此
\[
x^2+y^2=2y
\]
即曲线 \(C\) 的直角坐标方程为 \(x^2+y^2-2y=0\).
\end{solution}

\begin{problem}
若直线 \(3x - 4y + 5 = 0\) 与圆 \(x^{2} + y^{2} = r^{2}\) （\(r > 0\)） 相交于 \(A\)，\(B\) 两点，且 \(\angle AOB = 120^{\circ}\) （\(O\) 为坐标原点），则 \(r = \)\fillinblank{}.
\end{problem}

\begin{answer}
\(2\)
\end{answer}

\begin{solution}
圆心 \(O\) 到直线 \(3x-4y+5=0\) 的距离为
\[
d=\frac{|5|}{\sqrt{3^2+(-4)^2}}=1
\]
圆心到弦 \(AB\) 的垂线平分弦 \(AB\)，设垂足为 \(M\). 在直角三角形 \(OMA\) 中，
\[
\angle AOM=\frac12\angle AOB=60^{\circ}
\]
所以
\[
d=OM=r\cos60^{\circ}=\frac r2
\]
于是 \(r=2\).
\end{solution}

\begin{problem}
若函数 \( f(x) = |2^x - 2| - b \) 有两个零点，则实数 \( b \) 的取值范围是\fillinblank{}.
\end{problem}

\begin{answer}
\((0,2)\)
\end{answer}

\begin{solution}
函数有两个零点等价于方程
\[
|2^x-2|=b
\]
有两个不同的实数解. 必须先有 \(b>0\). 分情况去掉绝对值，得到
\[
2^x=2+b
\quad\text{或}\quad
2^x=2-b
\]
第一式对任意 \(b>0\) 有唯一解；第二式有实数解当且仅当 \(2-b>0\)，即 \(b<2\). 当 \(0<b<2\) 时两式右端不同且均为正，因而恰有两个不同的解；\(b=0\) 时两式给出同一个解，\(b\geqslant2\) 时第二式无实数解. 所以 \(b\) 的取值范围为 \((0,2)\).
\end{solution}

\begin{problem}
已知 \(\omega > 0\)，在函数 \(y = 2\sin \omega x\) 与 \(y = 2\cos \omega x\) 的图象的交点中，距离最短的两个交点的距离为 \(2\sqrt{3}\)，则 \(\omega = \)\fillinblank{}.
\end{problem}

\begin{answer}
\(\frac{\pi}{2}\)
\end{answer}

\begin{solution}
两图象的交点满足
\[
2\sin\omega x=2\cos\omega x
\]
即 \(\tan\omega x=1\). 因而交点可记为
\(P_k=\left(\frac{\pi/4+k\pi}{\omega},\,\sqrt2(-1)^k\right)\)，\(k\in\Z\).
令 \(t=\dfrac\pi\omega>0\). 相邻两个交点的横坐标相差 \(t\)，纵坐标相差 \(2\sqrt2\)，所以相邻交点的距离为
\[
d_1=\sqrt{t^2+8}
\]
横坐标相差两个间隔的两个交点纵坐标相同，距离为
\[
d_2=2t
\]
横坐标相差更多间隔的点对距离不小于上述两类距离，因此最短距离为 \(\min\{d_1,d_2\}\). 题设给出
\[
\min\{\sqrt{t^2+8},2t\}=2\sqrt3
\]
若 \(\sqrt{t^2+8}=2\sqrt3\)，则 \(t=2\). 若 \(2t=2\sqrt3\)，则 \(t=\sqrt3\)，但此时 \(\sqrt{t^2+8}=\sqrt{11}<2\sqrt3\)，与最短距离为 \(2\sqrt3\) 矛盾. 所以 \(t=2\)，从而
\[
\omega=\frac\pi t=\frac\pi2
\]
\end{solution}

\section{解答题}

\begin{problem}
某商场举行有奖促销活动，顾客购买一定金额的商品后即可抽奖，抽奖方法是：从装有 \(2\) 个红球 \(A_{1}\)，\(A_{2}\) 和 \(1\) 个白球 \(B\) 的甲箱与装有 \(2\) 个红球 \(a_{1}\)，\(a_{2}\) 和 \(2\) 个白球 \(b_{1}\)，\(b_{2}\) 的乙箱中，各随机摸出 \(1\) 个球. 若摸出的 \(2\) 个球都是红球则中奖，否则不中奖.
\begin{enumerate}
\item 用球的标号列出所有可能的摸出结果；
\item 有人认为：两个箱子中的红球比白球多，所以中奖的概率大于不中奖的概率，你认为正确吗？请说明理由.
\end{enumerate}
\end{problem}

\begin{solution}
\begin{enumerate}
\item 甲箱有 \(A_{1}\)，\(A_{2}\)，\(B\) 三种取法，乙箱有 \(a_{1}\)，\(a_{2}\)，\(b_{1}\)，\(b_{2}\) 四种取法. 由于两箱分别取球，所有可能的结果为
\[
\begin{gathered}
\{(A_{1},a_{1}),(A_{1},a_{2}),(A_{1},b_{1}),(A_{1},b_{2}),\\
(A_{2},a_{1}),(A_{2},a_{2}),(A_{2},b_{1}),(A_{2},b_{2}),\\
(B,a_{1}),(B,a_{2}),(B,b_{1}),(B,b_{2})\}.
\end{gathered}
\]
共 \(3\times4=12\) 种.

\item 每一种结果的概率都为 \(\dfrac{1}{12}\). 中奖的结果必须是甲箱取到红球且乙箱取到红球，即
\[
\{(A_{1},a_{1}),(A_{1},a_{2}),(A_{2},a_{1}),(A_{2},a_{2})\}
\]
因此中奖的概率为
\[
P(\text{中奖})=\frac{4}{12}=\frac{1}{3}
\]
不中奖是中奖的对立事件，所以
\[
P(\text{不中奖})=1-\frac{1}{3}=\frac{2}{3}
\]
因为 \(\dfrac{1}{3}<\dfrac{2}{3}\)，所以上述认为不正确.
\end{enumerate}
\end{solution}

\begin{problem}
设 \(\triangle ABC\) 的内角 \(A\)，\(B\)，\(C\) 的对边分别为 \(a\)，\(b\)，\(c\)，\(a = b \tan A\).
\begin{enumerate}
\item 证明：\(\sin B = \cos A\)；
\item 若 \(\sin C - \sin A\cos B = \frac{3}{4}\)，且 \(B\) 为钝角，求 \(A\)，\(B\)，\(C\).
\end{enumerate}
\end{problem}

\begin{solution}
\begin{enumerate}
\item 由正弦定理，存在同一个正数 \(2R\)，使得
\[
\frac{a}{\sin A}=\frac{b}{\sin B}=2R
\]
因此
\[
\frac{a}{b}=\frac{\sin A}{\sin B}
\]
由题设 \(a=b\tan A\)，且 \(a\)，\(b>0\)，知 \(\tan A>0\). 又三角形内角满足 \(0<A<\pi\)，所以 \(0<A<\dfrac{\pi}{2}\)，从而 \(\sin A>0\)、\(\cos A>0\)，并且
\[
\frac{a}{b}=\tan A=\frac{\sin A}{\cos A}
\]
比较两式可得 \(\dfrac{\sin A}{\sin B}=\dfrac{\sin A}{\cos A}\)，故
\[
\sin B=\cos A
\]

\item 因为 \(A+B+C=\pi\)，所以 \(C=\pi-(A+B)\)，进而
\[
\sin C=\sin(A+B)=\sin A\cos B+\cos A\sin B
\]
于是题设等式化为
\[
\sin C-\sin A\cos B=\cos A\sin B=\frac{3}{4}
\]
结合第（1）问的结论 \(\sin B=\cos A\)，得到
\[
\cos^{2}A=\frac{3}{4}
\]
由 \(\sin B=\cos A>0\)，知 \(\cos A>0\)，故 \(0<A<\dfrac{\pi}{2}\)，从而
\(\cos A=\frac{\sqrt{3}}{2}\)，\(A=\frac{\pi}{6}\).
此时 \(\sin B=\dfrac{\sqrt{3}}{2}\). 又因为 \(B\) 为钝角，所以
\[
B=\frac{2\pi}{3}
\]
最后
\[
C=\pi-A-B=\pi-\frac{\pi}{6}-\frac{2\pi}{3}=\frac{\pi}{6}
\]
故 \(A=\dfrac{\pi}{6}\)，\(B=\dfrac{2\pi}{3}\)，\(C=\dfrac{\pi}{6}\).
\end{enumerate}
\end{solution}

\begin{problem}
如图，直三棱柱 \(ABC - A_{1}B_{1}C_{1}\) 的底面是边长为 \(2\) 的正三角形，\(E\)，\(F\) 分别是 \(BC\)，\(CC_{1}\) 的中点.
\begin{enumerate}
\item 证明：平面 \(AEF \perp\) 平面 \(B_{1}BCC_{1}\)；
\item 若直线 \(A_{1}C\) 与平面 \(A_{1}ABB_{1}\) 所成的角为 \( 45^{\circ} \)，求三棱锥 \(F\text{-}AEC\) 的体积.
\end{enumerate}

\begin{minipage}{\linewidth}
    \centering
    \bitmapfigure[max width=0.42\linewidth]{img/2015/hunan_liberal/q18_fig1.png}
\end{minipage}
\end{problem}

\begin{solution}
\begin{enumerate}
\item 因为底面 \(ABC\) 是正三角形，且 \(E\) 是 \(BC\) 的中点，所以 \(AE\perp BC\). 又因为三棱柱是直棱柱，故 \(CC_{1}\perp\) 平面 \(ABC\)，从而 \(AE\perp CC_{1}\).

平面 \(B_{1}BCC_{1}\) 内的两条相交直线 \(BC\) 与 \(CC_{1}\) 分别与平面 \(AEF\) 内的直线 \(AE\) 垂直，因此 \(AE\perp\) 平面 \(B_{1}BCC_{1}\). 又因 \(AE\subset\) 平面 \(AEF\)，所以平面 \(AEF\perp\) 平面 \(B_{1}BCC_{1}\).

\item 设直三棱柱的高为 \(h\). 建立空间直角坐标系，取
\(B(-1,0,0)\)，\(C(1,0,0)\)，\(A(0,\sqrt{3},0)\)
使底面位于 \(xOy\) 平面内，则
\(A_{1}(0,\sqrt{3},h)\)，\(B_{1}(-1,0,h)\)，\(C_{1}(1,0,h)\).
平面 \(A_{1}ABB_{1}\) 是过直线 \(AB\) 且沿竖直方向的平面. 在底面内，直线 \(AB\) 的方程为 \(y=\sqrt{3}x+\sqrt{3}\)，所以平面 \(A_{1}ABB_{1}\) 的一个法向量可取
\[
\bs{n}=(\sqrt{3},-1,0)
\]
直线 \(A_{1}C\) 的一个方向向量为
\[
\bs{v}=\overrightarrow{A_{1}C}=(1,-\sqrt{3},-h)
\]
设直线 \(A_{1}C\) 与平面 \(A_{1}ABB_{1}\) 所成的角为 \(\theta\)，则
\[
\sin\theta=\frac{|\bs{v}\cdot\bs{n}|}{|\bs{v}|\,|\bs{n}|}
 =\frac{|\sqrt{3}+\sqrt{3}|}{\sqrt{4+h^{2}}\times2}
 =\frac{\sqrt{3}}{\sqrt{4+h^{2}}}
\]
由 \(\theta=45^{\circ}\)，得
\[
\frac{\sqrt{3}}{\sqrt{4+h^{2}}}=\frac{\sqrt{2}}{2}
\]
即 \(h^{2}=2\)，所以 \(h=\sqrt{2}\).

三角形 \(AEC\) 的底 \(EC=1\)，高 \(AE=\sqrt{3}\)，故
\[
S_{\triangle AEC}=\frac{1}{2}\times1\times\sqrt{3}=\frac{\sqrt{3}}{2}
\]
点 \(F\) 是 \(CC_{1}\) 的中点，所以它到平面 \(AEC\) 的距离为 \(\dfrac{h}{2}=\dfrac{\sqrt{2}}{2}\). 因此三棱锥 \(F\text{-}AEC\) 的体积为
\[
V=\frac{1}{3}S_{\triangle AEC}\cdot\frac{h}{2}
 =\frac{1}{3}\times\frac{\sqrt{3}}{2}\times\frac{\sqrt{2}}{2}
 =\frac{\sqrt{6}}{12}
\]
\end{enumerate}
\end{solution}

\begin{problem}
设数列 \(\{a_n\}\) 的前 \(n\) 项和为 \(S_{n}\). 已知 \(a_1 = 1\)，\(a_2 = 2\)，且 \(a_{n+2} = 3S_{n} - S_{n+1} + 3\)，\(n \in \N^*\).

\begin{enumerate}
\item 证明： \(a_{n+2} = 3a_n\)；
\item 求 \(S_{n}\).
\end{enumerate}
\end{problem}

\begin{solution}
\begin{enumerate}
\item 当 \(n=1\) 时，由已知条件
\[
a_{3}=3S_{1}-S_{2}+3=3a_{1}-(a_{1}+a_{2})+3=3-3+3=3=3a_{1}
\]
当 \(n\geqslant2\) 时，由递推式分别取 \(n\) 和 \(n-1\)，得
\(a_{n+2}=3S_{n}-S_{n+1}+3\)，\(a_{n+1}=3S_{n-1}-S_{n}+3\).
又因为 \(S_{n+1}=S_n+a_{n+1}\)，所以
\[
a_{n+2}=2S_n-a_{n+1}+3
\]
将第二个等式代入，得
\[
a_{n+2}=2S_n-(3S_{n-1}-S_n+3)+3
 =3S_n-3S_{n-1}=3(S_n-S_{n-1})=3a_n
\]
结合 \(n=1\) 时的结果，故对一切 \(n\in\N^{*}\)，都有
\[
a_{n+2}=3a_n
\]

\item 由 \(a_1=1\)、\(a_2=2\) 及上面的递推关系，奇数项和偶数项分别构成等比数列：
\(a_{2k-1}=3^{k-1}\)，\(a_{2k}=2\times3^{k-1}\)，\(k\in\N^{*}\).
若 \(n=2m\)，则
\[
S_{2m}=\sum_{k=1}^{m}a_{2k-1}+\sum_{k=1}^{m}a_{2k}
 =\sum_{k=1}^{m}3^{k-1}+2\sum_{k=1}^{m}3^{k-1}
 =3\frac{3^{m}-1}{3-1}
 =\frac{3(3^{m}-1)}{2}
\]
若 \(n=2m-1\)，则
\[
S_{2m-1}=\sum_{k=1}^{m}a_{2k-1}+\sum_{k=1}^{m-1}a_{2k}
 =\frac{3^{m}-1}{2}+2\frac{3^{m-1}-1}{3-1}
 =\frac{5\cdot3^{m-1}-3}{2}
\]
因此
\[
S_n=
\begin{cases}
\dfrac{5\cdot3^{m-1}-3}{2},&n=2m-1,\\[6pt]
\dfrac{3(3^m-1)}{2},&n=2m,
\end{cases}
\qquad m\in\N^{*}
\]
\end{enumerate}
\end{solution}

\begin{problem}
已知抛物线 \(C_1: x^2 = 4y\) 的焦点 \(F\) 也是椭圆 \(C_2: \frac{y^2}{a^2} + \frac{x^2}{b^2} = 1\)（\(a > b > 0\)） 的一个焦点，\(C_1\) 与 \(C_2\) 的公弦长为 \(2\sqrt{6}\). 过点 \(F\) 的直线 \(l\) 与 \(C_1\) 相交于 \(A\)，\(B\) 两点，与 \(C_2\) 相交于 \(C\)，\(D\) 两点，且 \(\overrightarrow{AC}\) 与 \(\overrightarrow{BD}\) 同向.
\begin{enumerate}
\item 求 \(C_{2}\) 的方程；
\item 若 \( \left|AC\right|=\left|BD\right| \)，求直线 l 的斜率.
\end{enumerate}
\end{problem}

\begin{solution}
\begin{enumerate}
\item 由 \(C_{1}:x^{2}=4y\) 可知其焦点为 \(F(0,1)\). 由于 \(F\) 也是椭圆 \(C_{2}\) 的焦点，且椭圆的长轴在 \(y\) 轴上，因此
\[
a^{2}-b^{2}=1
\]
将抛物线方程中的 \(x^{2}=4y\) 代入椭圆方程，公共点的纵坐标 \(y\) 满足
\[
\frac{y^{2}}{a^{2}}+\frac{4y}{b^{2}}=1
\]
对 \(y>0\)，左边关于 \(y\) 严格递增，所以抛物线与椭圆只有一条水平公共弦. 设其端点纵坐标为 \(y_{0}\)，则公共弦端点为 \((\pm2\sqrt{y_{0}},y_{0})\)，弦长为 \(4\sqrt{y_{0}}\). 由题设
\[
4\sqrt{y_{0}}=2\sqrt{6}
\]
得 \(y_{0}=\dfrac{3}{2}\).

代入上式并令 \(u=a^{2}\)，再利用 \(b^{2}=a^{2}-1=u-1\)，有
\[
\frac{9}{4u}+\frac{6}{u-1}=1
\]
因为 \(a>b>0\)，所以 \(u>1\). 化简得
\[
4u^{2}-37u+9=0
\]
故
\[
u=\frac{37\pm35}{8}
\]
舍去不满足 \(u>1\) 的值 \(\dfrac14\)，得到 \(a^{2}=9\)，进而 \(b^{2}=8\). 因此
\[
C_{2}:\frac{y^{2}}{9}+\frac{x^{2}}{8}=1
\]

\item 设直线 \(l\) 的斜率为 \(k\)，则其方程为
\[
y=kx+1
\]
记 \(s=\sqrt{k^{2}+1}\). 将直线方程代入抛物线，得
\[
x^{2}-4kx-4=0
\]
所以抛物线交点 \(A\)，\(B\) 的横坐标可按从小到大记为
\(x_{A}=2(k-s)\)，\(x_{B}=2(k+s)\)
从而
\[
x_{B}-x_{A}=4s
\]

将直线方程代入椭圆方程，得
\[
\frac{x^{2}}{8}+\frac{(kx+1)^{2}}{9}=1
\]
即
\[
(9+8k^{2})x^{2}+16kx-64=0
\]
其两根按从小到大记为 \(x_{C}\)，\(x_{D}\). 由根的差值公式，判别式为
\[
\Delta=(16k)^{2}+4(9+8k^{2})\times64
 =256\left(k^{2}+9+8k^{2}\right)
 =2304s^{2}
\]
所以
\[
x_{D}-x_{C}=\frac{\sqrt{\Delta}}{9+8k^{2}}
 =\frac{48s}{9+8k^{2}}
\]

直线上的两点横坐标之差乘以 \(s\) 就是两点间距离. 又因 \(\overrightarrow{AC}\) 与 \(\overrightarrow{BD}\) 同向且 \(\left|AC\right|=\left|BD\right|\)，两段在直线上的有向位移相等，故
\[
x_{C}-x_{A}=x_{D}-x_{B}
\]
移项可得
\[
x_{B}-x_{A}=x_{D}-x_{C}
\]
代入上面得到的两条弦长关系，得
\[
4s=\frac{48s}{9+8k^{2}}
\]
由于 \(s>0\)，所以 \(9+8k^{2}=12\)，即
\[
k^{2}=\frac{3}{8}
\]
因此直线 \(l\) 的斜率为
\[
k=\pm\frac{\sqrt{6}}{4}
\]
\end{enumerate}
\end{solution}

\begin{problem}
已知 \(a > 0\)，函数 \(f(x) = a\e^{x}\cos x\)，\(x \in [0, +\infty)\). 记 \(x_n\) 为 \(f(x)\) 的从小到大的第 \(n\) 个极值点（\(n \in \N^*\)）.
\begin{enumerate}
\item 证明：数列 \( \{f(x_{n})\} \) 是等比数列；
\item 若对一切 \(n \in \N^*\)，\(x_n \leqslant |f(x_n)|\) 恒成立，求 \(a\) 的取值范围.
\end{enumerate}
\end{problem}

\begin{solution}
\begin{enumerate}
\item 对函数求导，得
\[
f'(x)=a\e^{x}(\cos x-\sin x)
\]
由于 \(a\e^{x}>0\)，且
\[
\cos x-\sin x=\sqrt{2}\cos\left(x+\frac{\pi}{4}\right)
\]
所以导函数的零点为
\(x=\frac{\pi}{4}+k\pi\)，\(k\in\Z\).
该因子在每个零点两侧变号，故这些零点都是极值点，且没有其他极值点. 在 \([0,+\infty)\) 内，按从小到大的顺序可得
\(x_{n}=\frac{\pi}{4}+(n-1)\pi\)，\(n\in\N^{*}\).
于是
\[
f(x_{n})=a\e^{x_{n}}\cos\left(\frac{\pi}{4}+(n-1)\pi\right)
 =\frac{a\e^{\frac{\pi}{4}+(n-1)\pi}}{\sqrt{2}}(-1)^{n-1}
\]
因此
\[
\frac{f(x_{n+1})}{f(x_{n})}=-\e^{\pi}
\]
它与 \(n\) 无关，所以数列 \(\{f(x_{n})\}\) 是等比数列，首项为 \(\dfrac{a\e^{\pi/4}}{\sqrt{2}}\)，公比为 \(-\e^{\pi}\).

\item 由上式
\[
|f(x_{n})|=\frac{a\e^{x_{n}}}{\sqrt{2}}
\]
题设不等式等价于
\(a\geqslant\sqrt{2}x_{n}\e^{-x_{n}}\)，\(n\in\N^{*}\).
令 \(g(x)=x\e^{-x}\)，则
\[
g'(x)=(1-x)\e^{-x}
\]
所以 \(g(x)\) 在 \((0,1)\) 上递增，在 \((1,+\infty)\) 上递减. 其中
\(x_{1}=\frac{\pi}{4}<1\)，\(x_{2}=\frac{5\pi}{4}>1\).
对 \(n\geqslant2\)，有 \(x_n\geqslant x_2\)，故 \(g(x_n)\leqslant g(x_2)\). 又
\[
\frac{g(x_{1})}{g(x_{2})}
 =\frac{x_{1}}{x_{2}}\e^{x_{2}-x_{1}}
 =\frac{\e^{\pi}}{5}>1
\]
其中 \(\e^{\pi}>1+\pi+\dfrac{\pi^{2}}{2}+\dfrac{\pi^{3}}{6}>13>5\)，所以该比较成立.
所以 \(g(x_{1})>g(x_{2})\)，从而所有 \(g(x_n)\) 中的最大值为 \(g(x_1)\). 因此
\[
a\geqslant\sqrt{2}g(x_{1})
 =\sqrt{2}\times\frac{\pi}{4}\e^{-\frac{\pi}{4}}
 =\frac{\pi}{2\sqrt{2}}\e^{-\frac{\pi}{4}}
\]
故 \(a\) 的取值范围为
\[
\left[\frac{\pi}{2\sqrt{2}}\e^{-\frac{\pi}{4}},+\infty\right)
\]
\end{enumerate}
\end{solution}
